\scanpage{78}
\uline{Remarque} Les cas I, II \inserted{\deleted{\unsure{et}} II'} correspondent \inserted{chacun} à une structure, définie à isom. (unique près), \inserted{qui existe \deleted{pour} $\forall\,k$ de car $0$ (i.e. définie sur $\mathbb Q$)} les autres \inserted{trois} correspondent chacun à deux structures, définies à isom unique près, \deleted{celles} qui \unsure{existe|existent} ssi $k$ contient une solution de $\alpha^2+\alpha-1=0$, i.e. ssi $k$ contient le corps $\tfrac12$ cyclotomique d'indice $5$. Il y a donc en tout $3+3\times2=$\uuline{$9$} types d'isomorphie de ``polyèdres réguliers finis'' sur corps $k$ de car $0$, contenant $\alpha$.

\uline{\unsure{Indication|Indications} sur dém. du th.} \deleted{\unsure{Lisse :|Lisons}} OPS \inserted{$k$ alg. clos, et $=$} $k=\mathbb C$.

I) Notons d'abord que d'après la prop, si valeurs $\nu_0,\nu_1$ sont finies, on a \deleted{\unsure{$\beta_0\beta_1$ inv}} $\alpha_0\ne\pm2,\alpha_1\ne\pm2$

\deleted{\unsure{d'où} $\sigma,\sigma'\ne0$, donc $Q_X$ lisse. Peut-on avoir $\rho=0$ i.e. $\alpha_0+\alpha_1=0$ ? i.e. $\alpha_1=-2\alpha_0$ \quad $\alpha_0=\alpha,\alpha_1=-\alpha$, où $\alpha$ est une solution d'une équation $\tfrac12$ cyclotomique (donnée par un indice $\nu=\nu_0\geq\unsure{3|2|3/2}$ ? Comme \unsure{$\rho\ne0$} $\sigma=2-\alpha_1=2\alpha\ne0$, on serait dans un cas où $Q_E=Q\setminus Q\cap C=Q\setminus Y$ est un plan \unsure{euclidien}}

Les $\alpha_i$ sont solutions d'équations $\tfrac12$-cyclotomiques d'indices $\nu_0,\nu_1\geq3$. Ces solutions sont \uline{réelles}, donc la situation provient d'une situation sur $k_0=\mathbb R$. Comme on prend $G_0$ \deleted{\unsure{est|étant}} fini, $\exists$ forme quadratique invariante définie positive sur $V$, \deleted{\unsure{et sur $E$ donc}} $\boxed{\beta\rho\ne0}$ i.e. \uline{$Y$ lisse}.

$\unsure{\mathbb G_0|G_0|\underline G_0}\subset\operatorname{GP}(1,\mathbb C)$ groupe fini d'autom de $Y\simeq\mathbb P^1_{\mathbb C}$, \deleted{\unsure{utilisons}} soit $Z=Y/G_0$ \inserted{$\simeq\mathbb P^1_{\mathbb C}$} \unsure{Utilisant|Utilisons} \unsure{Hurwicz|Hurwitz} on trouve que $Y$ est ramifié sur $Z$ en $2$ ou trois points \inserted{exactement}. Dans le cas de $2$ pts, on aurait au dessus \unsure{des|de} ceux-ci un diviseur de degré $2$ stable sous $G_0$, absurde. Donc il y a trois points de ramification, et la formule de Hurwicz nous dit que les \deleted{\unsure{indices de}} indices de ramification sont $(2,p,q)$ -- les couples $(p,q)$ sont \inserted{parmi} ceux inventoriés plus haut, \inserted{savoir} $(3,3)(3,4)(3,5)$ \deleted{(et} (à \unsure{la} symétrie près), plus $(2,\unsure{2|p},p)$ (cas diédral) -- mais dans celui-ci, il y a au dessus des

\margin{$\mathbb G_0=\operatorname{Im}G_0$, $G_0\to\mathbb G_0$ de noyau d'ordre 1 ou 2.}
% archive sheet number: 77
